\documentclass[12pt]{article}
\usepackage{latexblog}

% ============================================================
% Blog Metadata
% ============================================================
\blogtitle{理解Java并发(1)：基本机制}
\blogdate{2019-10-28}
\blogtags{刨根问底, ~理解Java并发, 并发编程, Concurrent}
\bloglang{zh}
\blogsummary{}

\begin{document}
\makeblogtitle

线程是操作系统中进行运算调度的最小单位，它是一个单一顺序的控制流，不论是对于单核还是多核的CPU，都能比较有效的提高程序的吞吐率。在Java中，创建一个线程的唯一方法是创建一个\texttt{Thread}类的实例，并调用\texttt{start()}方法以启动该线程。然而当多个线程同时执行时，如何保证线程之间是按照我们期待的方式在运行呢？Java提供了多种机制来保证多个线程之间的交互。

\section{同步(Synchronization)与监视器(Monitor)机制}\label{ux540cux6b65synchronizationux4e0eux76d1ux89c6ux5668monitorux673aux5236}

显而易见最基本最常见的和多线程有关的就是同步\texttt{synchronized}关键字了，它底层是使用Monitor实现的。那么究竟什么是\texttt{Monitor}呢？根据JavaSE Specification的描述，在Java中，每一个对象都有一个与之关联的monitor，允许线程可以去\texttt{lock}或者\texttt{unlock}这个monitor。实际上：

\begin{itemize}
\tightlist
\item
  \texttt{monitor}是独立于Java语言之上的一个概念（没想到还有另外一个名字\texttt{管程}），保证在运行线程之前获取互斥锁
\item
  在Java中，任何对象(\texttt{java.lang.Object})都可以允许作为一个monitor，所以会有\texttt{wait}、\texttt{notify}之类的方法
\end{itemize}

\texttt{synchronized}可以作用于代码块或者方法上。如果作用在代码块上，它会尝试去lock这个对象的monitor，如果不成功将会等待直到lock成功。而当执行完毕后，无论是否出现异常，都将会释放这个锁。

如果作用在方法上，唯一的区别在于，如果是实例方法，那么将使用这个实例作为monitor，也就是\texttt{this}；如果是静态方法，那么使用的是所在类的\texttt{Class}对象。

\section{Wait/Notify}\label{waitnotify}

每一个Object都包含一个等待线程的集合(Wait set)。当对象创建的时候，这个队列是空的，当调用\texttt{Object.wait()}、\texttt{Object.nofity()}以及\texttt{Object.nofityAll()}方法的时候，会自动添加或者移除队列中的线程。或者当线程的中断状态发生改变的时候，也会引起变化。

\subsection{Wait}\label{wait}

调用\texttt{wait}方法将使当前线程休眠直到另一个线程通过\texttt{notify}或者\texttt{notifyAll}来唤醒。当前线程必须持有该对象的锁，调用\texttt{wait}后即释放锁。当线程被唤醒时，需要重新取得锁并继续执行。然而，线程被唤醒有可能是因为``虚假唤醒''（spurious wakeups）导致，所以通常都需要将\texttt{wait}检测的逻辑包括在一个loop中：

\begin{Shaded}
\begin{Highlighting}[]
\KeywordTok{synchronized} \OperatorTok{(}\NormalTok{obj}\OperatorTok{)} \OperatorTok{\{}
    \ControlFlowTok{while} \OperatorTok{(\textless{}}\NormalTok{condition does not hold}\OperatorTok{\textgreater{})}
\NormalTok{        obj}\OperatorTok{.}\FunctionTok{wait}\OperatorTok{();}
    \CommentTok{// Perform action appropriate to condition}
\OperatorTok{\}}
\end{Highlighting}
\end{Shaded}

所谓虚假唤醒就是说，本来不该唤醒的时候唤醒了。究其原因是在操作系统层面就性能和正确性做出了权衡，放弃了正确性而选择让程序自己去处理。

\begin{quote}
Spurious wakeups may sound strange, but on some multiprocessor systems, making condition wakeup completely predictable might substantially slow all condition variable operations.
\end{quote}

\subsection{Notify}\label{notify}

调用\texttt{notify}将唤醒一个正在等待持有该对象锁的线程，如果有多个对象在等待的话，将会随机唤醒其中的一个。

被唤醒的线程必须等到当前线程释放锁之后，才能开始执行；也就是说\texttt{notify}执行完之后，并不会立即释放锁，而是需要等到同步块执行完。

如果调用\texttt{notifyAll}的话，所有等待的线程将被唤醒，但同一时间有且仅有一个线程能取到锁并继续执行。

\subsection{Interruption}\label{interruption}

当调用\texttt{Thread.interrupt}时，线程的中断状态呗设置为true。如果该线程在某个对象的waitSet中，则将会被从等待队列中移除，并在取得锁之后抛出\texttt{InterruptedException}。实际上，如果线程正在执行的是一些底层的blocking函数例如\texttt{Thread.sleep()}, \texttt{Thread.join()}, 或者 \texttt{Object.wait()}的时候，那么线程将抛出\texttt{InterruptedException}，并且\texttt{interrupted}状态会被清除；否则只会将\texttt{interrupted}状态设置为\texttt{true}。

如果一个处于等待队列中的线程同时收到中断和通知，那么可能的行为是：

\begin{itemize}
\tightlist
\item
  先收到通知，正常唤醒。这时候，\texttt{Thread.interrupted}将为\texttt{true}，
\item
  抛出\texttt{InterruptedException}并退出
\end{itemize}

同样，如果有多个线程处于对象m的等待队列中，然后另一个线程执行\texttt{m.notify}，那么可能：

\begin{itemize}
\tightlist
\item
  至少有一个线程正常退出wait
\item
  所有处于等待队列中的线程抛出\texttt{InterruptedException}而退出
\end{itemize}

需要注意的是，当一个线程中断了另一个线程的时候，被中断的线程并不是需要立即停止执行，程序可以选择在停止之前做一些清理工作之类的。通常如果捕获了\texttt{InterruptedException}只需要重新抛出即可，有些时候不能重新抛出的时候，需要将当前线程标记为\texttt{interrupted}使得上层堆栈的程序可以选择处理，

\begin{Shaded}
\begin{Highlighting}[]
\ControlFlowTok{try} \OperatorTok{\{}
    \ControlFlowTok{while} \OperatorTok{(}\KeywordTok{true}\OperatorTok{)} \OperatorTok{\{}
\NormalTok{        Task task }\OperatorTok{=}\NormalTok{ queue}\OperatorTok{.}\FunctionTok{take}\OperatorTok{(}\DecValTok{10}\OperatorTok{,} \BuiltInTok{TimeUnit}\OperatorTok{.}\FunctionTok{SECONDS}\OperatorTok{);}
\NormalTok{        task}\OperatorTok{.}\FunctionTok{execute}\OperatorTok{();}
    \OperatorTok{\}}
\OperatorTok{\}}\ControlFlowTok{catch} \OperatorTok{(}\BuiltInTok{InterruptedException}\NormalTok{ e}\OperatorTok{)} \OperatorTok{\{} 
    \BuiltInTok{Thread}\OperatorTok{.}\FunctionTok{currentThread}\OperatorTok{().}\FunctionTok{interrupt}\OperatorTok{();}
\OperatorTok{\}}
\end{Highlighting}
\end{Shaded}

\section{线程的生命周期}\label{ux7ebfux7a0bux7684ux751fux547dux5468ux671f}

每一个线程有一个生命周期，包含多个状态：

\begin{itemize}
\tightlist
\item
  New: 线程创建还未开始执行，线程创建完之后即为此状态
\item
  Runnable: 在JVM中正在执行的状态。当线程start之后，即变为runnable状态
\item
  Blocked: 线程等待获取锁而被阻塞
\item
  Waiting: 线程等待其他线程
\item
  Timed Waiting: 有超时的等待
\item
  Terminated: 线程已被退出
\end{itemize}

\begin{figure}
\centering
\pandocbounded{\includegraphics[keepaspectratio,alt={Life cycle of a thread}]{https://www.baeldung.com/wp-content/uploads/2018/02/Life_cycle_of_a_Thread_in_Java.jpg}}
\caption{Life cycle of a thread}
\end{figure}

\section{Sleep / Yield}\label{sleep-yield}

当调用线程的\texttt{sleep}方法将导致线程暂时停止执行，值得注意的是并不会释放锁。而当线程的\texttt{yield}方法被调用时，意味着通知CPU当前线程可以``暂缓''执行的，实际很少使用。

\begin{quote}
It is rarely appropriate to use this method. It may be useful for debugging or testing purposes, where it may help to reproduce bugs due to race conditions.
\end{quote}

\section{Context switching}\label{context-switching}

在多线程中，CPU会为每个线程分配时间片区执行，即执行当前线程的一部分操作之后，操作系统需要从当前线程切换到其他线程中去。通常在下列的情况下会出现context switching:

\begin{itemize}
\tightlist
\item
  多任务处理（即多个线程正常执行）
\item
  中断，
\end{itemize}

那么在这个切换的过程中，会发生一些什么事情呢？

参考：

\begin{itemize}
\tightlist
\item
  \href{https://docs.oracle.com/javase/specs/jls/se7/html/jls-17.html}{Chapter 17. Threads and Locks}
\item
  \href{https://stackoverflow.com/questions/3362303/whats-a-monitor-in-java}{What's a monitor in Java?}
\item
  \href{https://zh.wikipedia.org/wiki/\%E7\%9B\%A3\%E8\%A6\%96\%E5\%99\%A8_(\%E7\%A8\%8B\%E5\%BA\%8F\%E5\%90\%8C\%E6\%AD\%A5\%E5\%8C\%96)}{管程}
\item
  \href{https://stackoverflow.com/questions/1050592/do-spurious-wakeups-in-java-actually-happen}{Do spurious wakeups in Java actually happen?}
\item
  \href{https://stackoverflow.com/questions/8594591/why-does-pthread-cond-wait-have-spurious-wakeups}{Why does pthread\_cond\_wait have spurious wakeups?}
\item
  \href{https://www.ibm.com/developerworks/java/library/j-jtp05236/index.html}{Dealing with InterruptedException}
\item
  \href{https://docs.oracle.com/javase/7/docs/api/java/lang/Thread.State.html}{Enum Thread.State}
\item
  \href{https://www.baeldung.com/java-thread-lifecycle}{Life Cycle of a Thread in Java}
\item
  \href{https://en.wikipedia.org/wiki/Context_switch}{Context switch}
\item
  \href{https://blog.tsunanet.net/2010/11/how-long-does-it-take-to-make-context.html}{How long does it take to make a context switch?}
\end{itemize}


\end{document}
